\lsection{前順序}

\lsubsection{薄い圏}

\dfn{薄い圏}{
  圏$\bm{\Lambda}$が以下を満たすとき、圏$\bm{\Lambda}$を薄いという。
  \eq*{
    \forall a, b \in \Obj_{\bm{\Lambda}} \forall f, g \in \Hom_{\bm{\Lambda}}\qty(a, b) \qty(f = g)
  }
}

\lem{薄い圏の簡約}{
  薄い圏$\bm{\Lambda}$について、任意の射は左簡約可能かつ右簡約可能である。
}{
  射$f, g, h \in \bm{\Lambda}$について、$f \circ g = f \circ h$とする。

  薄い圏であるので、$g = h$である。

  右簡約可能性も同様に示せる。
}

\dfn{大小}{
  薄い圏$\bm{\Lambda}$について、述語$\preccurlyeq_{\bm{\Lambda}}$を以下で定める。
  \eq*{
    a \preccurlyeq_{\bm{\Lambda}} b \defiff a \in \Obj_{\bm{\Lambda}} \land b \in \Obj_{\bm{\Lambda}} \land \exists f \in \Hom_{\bm{\Lambda}}\qty(a, b)
  }
}

\dfn{前順序に関わるいくつかの述語}{
  薄い圏$\bm{\Lambda}$について、述語$\prec, \succcurlyeq, \succ$を以下のように定める。
  \eqg*{
    a \prec b \defiff a \preccurlyeq b \land a \neq b \\*
    a \succcurlyeq b \defiff b \preccurlyeq a \\*
    a \succ b \defiff b \prec a
  }
}

\cor*{
  薄い圏$\bm{\Lambda}$について、以下が成り立つ。
  \eqg*{
    \forall a, b \in \bm{\Lambda} \qty(a \preccurlyeq b \leftrightarrow a \prec b \lor a = b) \\*
    \forall a, b \in \bm{\Lambda} \qty(a \succcurlyeq b \leftrightarrow a \succ b \lor a = b)
  }
}

\dfn{最大元}{
  薄い圏$\bm{\Lambda}$について、その終対象を最大元と呼ぶ。
}

\dfn{最小元}{
  薄い圏$\bm{\Lambda}$について、その始対象を最小元と呼ぶ。
}

\dfn{下限}{
  薄い圏$\bm{\Lambda}$、離散圏$\bm{J}$と、関手$X \colon \bm{J} \to \bm{\Lambda}$を考える。

  $X$の積$a, u$について、$a$を$X$の下限と呼ぶ。
}

\dfn{上限}{
  薄い圏$\bm{\Lambda}$、離散圏$\bm{J}$と、関手$X \colon \bm{J} \to \bm{\Lambda}$を考える。

  $X$の和$a, u$について、$a$を$X$の上限と呼ぶ。
}

\dfn{射影系}{
  薄い圏$\bm{\Lambda}$と圏$\bm{C}$について、関手$X \colon \bm{\Lambda} \to \bm{C}$を、射影系と呼ぶ。
}

\dfn{帰納極限}{
  射影系$X$ついて、$X$の極限を帰納極限と呼ぶ。
}

\dfn{射影極限}{
  射影系$X$ついて、$X$の余極限を射影極限と呼ぶ。
}


\lsubsection{前順序集合}

\rem{射としての前順序}{
  集合$P$と以下を満たす集合$\Lambda \subseteq P \times P$を考える。
  \eqg*{
    \forall a \in P \qty(\qty(a, a) \in \Lambda) \\*
    \forall a, b, c \in P \qty\big(\qty(a, b) \in \Lambda \land \qty(b, c) \in \Lambda \rightarrow \qty(a, c) \in \Lambda)
  }

  $\Lambda$の元について、$=, \dom, \cod, \Comp$を以下で定めることにより、\subsecref{射の公理}を満たす。

  $\qty(a, b), \qty(b, c), \qty(a^\prime, c^\prime) \in \Lambda$について、
  \begin{enumerate}
    \item 等号は集合の相等として定める。
    \item $\dom(\qty(a, b)) \coloneqq \qty(a, a)$
    \item $\cod(\qty(a, b)) \coloneqq \qty(b, b)$
    \item $\Comp(\qty(a, b), \qty(b, c), \qty(a^\prime, c^\prime)) \defiff \qty(a, c) = \qty(a^\prime, c^\prime)$
  \end{enumerate}

  このとき、$\Lambda$は薄い圏となる。
}

\dfn{前順序集合}{
  集合$P$と以下を満たす集合$\Lambda \subseteq P \times P$を考える。
  \eqg*{
    \forall a \in P \qty(\qty(a, a) \in \Lambda) \\*
    \forall a, b, c \in P \qty\big(\qty(a, b) \in \Lambda \land \qty(b, c) \in \Lambda \rightarrow \qty(a, c) \in \Lambda)
  }

  薄い圏$\Lambda$の大小$\preccurlyeq_\Lambda$を用いて、前順序集合の大小$\preccurlyeq_P$を以下で定める。
  \eq*{
    a \preccurlyeq_P b \defiff \qty(a, a) \preccurlyeq_\Lambda \qty(b, b)
  }

  このとき、$\qty(P, \preccurlyeq_P)$を前順序集合と呼ぶ。
}

\dfn{上界}{
  前順序集合$\qty(P, \preccurlyeq)$について、$P$の部分$A$を考える。

  $A$が上に有界であるとは、以下を満たすことである。
  \eq*{
    \exists b \in P \forall a \in A \qty(a \preccurlyeq b)
  }

  このとき、$b$を$A$の上界と呼ぶ。
}

\dfn{下界}{
  前順序集合$\qty(P, \preccurlyeq)$について、$P$の部分$A$を考える。

  $A$が下に有界であるとは、以下を満たすことである。
  \eq*{
    \exists b \in P \forall a \in A \qty(b \preccurlyeq a)
  }

  このとき、$b$を$A$の下界と呼ぶ。
}

\dfn{有界}{
  前順序集合$\qty(P, \preccurlyeq)$について、$P$の部分$A$を考える。

  $A$が有界であるとは、$A$が上に有界かつ下に有界であることである。
}

\cor*{
  空でない前順序集合$\qty(P, \preccurlyeq)$について、$\varnothing$は有界である。
}

単調、前順序の圏

\lsubsection{有向集合}

\lsubsection{同値関係}
